% SPDX-License-Identifier: CC-BY-SA-4.0
% Author: Matthieu Perrin
% Part: <Nom de la partie>
% Section: <Nom de la section>
% Sub-section: <Nom de la sous-section>  % (facultatif, laisser vide si non utilisé)
% Frame: <Titre de la slide>

\begingroup

\SetKwData{X}{x}
\SetKwData{I}{i}
\SetKwFunction{AlgoCompter}{compter}
\SetKwFunction{AlgoBoucler}{boucler}

\begin{frame}{Problème de l'arrêt}
  \small
 
  \begin{block}{Définition -- Le langage $L_{halt}$}
    Étant données une MTD $M$ et une entrée $u$, est-ce que $M$ termine sur $u$ ? 
 
    $$L_{\mathit{halt}} \eqdef \{ \langle \mathit{encode}(M), u \rangle \in (\{0,1\}^\star)^2  \mid M \text{ termine sur } u \}$$ 
    
  \end{block}
 
  \pause 
  \begin{block}{Propriétés}
    \begin{enumerate}
    \item<2-> $L_{\mathit{halt}}$ est \structure{récursivement énumérable}.
    \item<3-> $L_{\mathit{halt}}$ est \structure{indécidable}.
    \item<4-> $\overline{L_{\mathit{halt}}}$ n'est \structure{pas récursivement énumérable}.
    \end{enumerate}
  \end{block}
 
  \begin{block}{Démonstration}
    \begin{enumerate}
    \item<2-> $L_{\mathit{halt}}$ est reconnu par la MTD universelle normalisée pour la reconnaissance.
    \item<3-> Un peu plus difficile...
    \item<4-> Sinon, $L_{\mathit{halt}}$ et $\overline{L_{\mathit{halt}}}$ seraient récursivement énumérable, donc décidables.
    \end{enumerate}
  \end{block}
\end{frame}

\endgroup
